\section{Práctica con potencias y raíces}

En los siguientes ejercicios conviene indicar las propiedades utilizadas y comprobar que cada expresión esté definida antes de transformarla. Evaluar no siempre significa obtener un decimal: una potencia o un radical puede ser la representación exacta más clara del número.

\subsection*{Traducir productos y potencias}

Reescriba cada producto como una potencia. Luego evalúe cuando resulte sencillo.

\begin{enumerate}
\item \(2\cdot2\cdot2\cdot2\cdot2\)
\item \(7\cdot7\cdot7\)
\item \((-3)\cdot(-3)\cdot(-3)\cdot(-3)\)
\item \(\dfrac25\cdot\dfrac25\cdot\dfrac25\)
\item \(10\cdot10\cdot10\cdot10\cdot10\cdot10\)
\end{enumerate}

Ahora reescriba cada potencia como un producto de factores iguales.

\begin{enumerate}
\item \(4^3\)
\item \(3^5\)
\item \((-2)^4\)
\item \(\left(\dfrac12\right)^3\)
\item \(a^6\)
\end{enumerate}

\subsection*{Base, exponente y valor}

En cada expresión indique la base y el exponente. Luego encuentre su valor.

\begin{enumerate}
\item \(2^6\)
\item \(5^3\)
\item \((-4)^2\)
\item \(-4^2\)
\item \((-1)^9\)
\item \(9^1\)
\item \(12^0\)
\item \(0^5\)
\end{enumerate}

Explique por qué no se pide evaluar \(0^0\) junto con los ejemplos anteriores.

\subsection*{Aplicar propiedades con exponentes naturales}

Reescriba cada expresión como una sola potencia cuando sea posible y luego evalúe.

\begin{enumerate}
\item \(2^3\cdot2^5\)
\item \(7^2\cdot7^4\)
\item \(\dfrac{5^7}{5^3}\)
\item \(\dfrac{3^6}{3^6}\)
\item \((2^3)^4\)
\item \((5^2)^3\)
\item \(2^4\cdot3^4\)
\item \(\left(\dfrac23\right)^3\)
\end{enumerate}

Indique por qué no puede aplicarse la regla de igual base a \(2^3\cdot3^4\).

\subsection*{Reconocer igualdades falsas}

Decida si cada igualdad es verdadera para todos los valores en que aparecen expresiones definidas. Si es falsa, proporcione un contraejemplo.

\begin{enumerate}
\item \(a^m\cdot a^n=a^{m+n}\)
\item \(a^m+a^n=a^{m+n}\)
\item \((a\cdot b)^n=a^n\cdot b^n\)
\item \((a+b)^2=a^2+b^2\)
\item \((a^m)^n=a^{mn}\)
\item \(a^n=n^a\)
\end{enumerate}

En las igualdades que considere verdaderas, indique las condiciones necesarias sobre bases y exponentes.

\subsection*{Jerarquía de operaciones}

Agregue paréntesis para mostrar la estructura indicada por la precedencia y luego evalúe.

\begin{enumerate}
\item \(3+2^4\)
\item \(2\cdot3^2+1\)
\item \(5-2^3\)
\item \(-3^2+10\)
\item \((-3)^2+10\)
\item \(2+4\cdot2^3\)
\item \((2+4)\cdot2^3\)
\item \((2+3)^2\)
\end{enumerate}

\subsection*{Exponentes enteros negativos}

Reescriba cada potencia mediante un exponente positivo y evalúe.

\begin{enumerate}
\item \(2^{-3}\)
\item \(5^{-2}\)
\item \(10^{-4}\)
\item \((-3)^{-3}\)
\item \(\left(\dfrac23\right)^{-2}\)
\item \(\left(-\dfrac45\right)^{-2}\)
\end{enumerate}

Simplifique luego las siguientes expresiones.

\begin{enumerate}
\item \(3^2\cdot3^{-5}\)
\item \(\dfrac{2^3}{2^7}\)
\item \(\dfrac{5^{-2}}{5^{-4}}\)
\item \((2^{-3})^2\)
\item \(4^0\cdot4^{-2}\)
\end{enumerate}

Explique por qué \(0^{-2}\) no está definido y por qué \(2^{-3}\) no es un número negativo.

\subsection*{Evaluar y reconocer raíces}

Evalúe cada raíz en \(\mathbb R\) cuando sea posible. Si no representa un número real, indíquelo.

\begin{enumerate}
\item \(\sqrt{49}\)
\item \(\sqrt{0}\)
\item \(\sqrt[3]{27}\)
\item \(\sqrt[3]{-64}\)
\item \(\sqrt[4]{16}\)
\item \(\sqrt[4]{-16}\)
\item \(\sqrt[5]{-32}\)
\item \(\sqrt{2}\)
\end{enumerate}

En el último caso, indique si el resultado es natural, entero, racional o irracional.

\subsection*{Raíz principal y ecuaciones}

Resuelva cada ecuación en \(\mathbb R\). Distinga las soluciones de la raíz principal correspondiente.

\begin{enumerate}
\item \(x^2=36\)
\item \(x^2=0\)
\item \(x^2=-9\)
\item \(x^3=64\)
\item \(x^3=-125\)
\item \(x^4=81\)
\end{enumerate}

Luego evalúe:
\[
\sqrt{(-5)^2},
\qquad
\sqrt[3]{(-5)^3},
\qquad
\sqrt{x^2}\text{ para }x=-7.
\]

\subsection*{Radicales y exponentes racionales}

Reescriba cada radical como una potencia con exponente racional.

\begin{enumerate}
\item \(\sqrt a\)
\item \(\sqrt[3]{a^2}\)
\item \(\sqrt[5]{x^3}\)
\item \(\dfrac1{\sqrt[4]{b}}\)
\item \(\sqrt[7]{y}\)
\end{enumerate}

Suponga positivas las letras que aparecen. Ahora reescriba como radicales:

\begin{enumerate}
\item \(a^{1/3}\)
\item \(a^{3/5}\)
\item \(x^{7/2}\)
\item \(b^{-2/3}\)
\item \(y^{-1/4}\)
\end{enumerate}

\subsection*{Evaluar potencias racionales}

Evalúe en \(\mathbb R\) cuando sea posible.

\begin{enumerate}
\item \(9^{1/2}\)
\item \(27^{1/3}\)
\item \(16^{3/4}\)
\item \(32^{2/5}\)
\item \(8^{-2/3}\)
\item \(81^{-1/2}\)
\item \((-8)^{1/3}\)
\item \((-8)^{2/3}\)
\item \((-16)^{1/4}\)
\item \(0^{3/5}\)
\item \(0^{-3/5}\)
\end{enumerate}

Antes de decidir sobre una base negativa, reduzca el exponente a su forma irreducible:
\[
(-8)^{2/6},
\qquad
(-32)^{6/15},
\qquad
(-16)^{2/8}.
\]

\subsection*{Propiedades con exponentes racionales}

Suponga que todas las letras representan números reales positivos. Simplifique y exprese el resultado mediante una sola potencia cuando sea posible.

\begin{enumerate}
\item \(a^{1/2}\cdot a^{3/2}\)
\item \(\dfrac{x^{7/3}}{x^{1/3}}\)
\item \((b^{2/5})^5\)
\item \(2^{3/4}\cdot2^{5/4}\)
\item \(\dfrac{3^{1/2}}{3^{5/2}}\)
\item \((a\cdot b)^{2/3}\)
\item \(\left(\dfrac xy\right)^{-1/2}\)
\item \(x^{-2/3}\cdot x^{5/3}\)
\end{enumerate}

Explique por qué la condición de positividad permite aplicar estas propiedades sin ambigüedades.

\subsection*{Preguntas conceptuales}

Responda con una explicación breve y, cuando resulte útil, con un ejemplo.

\begin{enumerate}
\item ¿Qué información aportan la base y el exponente de una potencia?
\item ¿Por qué \(a^0=1\) se define solamente para \(a\ne0\) en este capítulo?
\item ¿Qué relación hay entre un exponente negativo y el inverso multiplicativo?
\item ¿Por qué \(\sqrt{25}=5\), aunque \(x^2=25\) tenga dos soluciones?
\item ¿Por qué \(\sqrt{a^2}=\pm a\) y no siempre \(a\)? ¿Qué es la raíz principal?
\item ¿Qué necesidad conduce de \(\mathbb Q\) a \(\mathbb R\)?
\item ¿Por qué una raíz de índice par no admite un radicando negativo en \(\mathbb R\)?
\item ¿Por qué conviene definir primero las potencias racionales para bases positivas?
\item ¿Por qué debe reducirse \(p/q\) antes de estudiar una base negativa?
\item Dé un ejemplo que muestre que \(\sqrt{a+b}\ne\sqrt a+\sqrt b\) en general.
\end{enumerate}

\subsection*{Construir expresiones}

Construya una expresión que cumpla cada condición. En muchos casos existen varias respuestas correctas.

\begin{enumerate}
\item Una potencia de base natural y exponente mayor que \(1\) cuyo valor sea \(64\).
\item Dos potencias con bases y exponentes diferentes que representen el mismo número.
\item Una potencia de base negativa cuyo valor sea positivo.
\item Una potencia de base negativa cuyo valor sea negativo.
\item Una potencia con exponente negativo cuyo valor sea \(\frac1{25}\).
\item Una raíz de índice impar y radicando negativo cuyo valor sea entero.
\item Una ecuación \(x^n=a\) que tenga exactamente dos soluciones reales.
\item Una ecuación \(x^n=a\) que no tenga soluciones reales.
\item Una potencia racional de base positiva cuyo valor sea \(8\).
\item Una potencia racional de base negativa que esté definida en \(\mathbb R\).
\end{enumerate}
